\documentclass[10pt]{article}
\usepackage[T1]{fontenc}
\usepackage{amsmath,amssymb}
\usepackage{booktabs,longtable,array}
\usepackage[margin=0.82in]{geometry}
\usepackage{enumitem}
\usepackage{xurl}
\usepackage[hidelinks]{hyperref}

\newcommand{\card}[1]{\lvert#1\rvert}
\newcommand{\Z}{\mathbb Z}
\newcolumntype{L}[1]{>{\raggedright\arraybackslash}p{#1}}
\setlist{leftmargin=*}
\renewcommand{\arraystretch}{1.08}

\title{Function Types Applicable to the Scheme in the Paper\\
A New Framework for Efficient Multivariate Functional Bootstrapping}
\author{}
\date{}

\begin{document}
\maketitle

\section{Why this note exists}

This note records which finite functions have actually compressed well under the
ideas used in the paper.  It is not a theorem claiming that every member of a
named family will compress.  We tried more functions than are shown in the main
tables, kept the successful parameter sets, and left all rejected rows in the raw
CSV and JSON files.

The common opportunity is simple:
\[
  \card{\operatorname{Im}f}\ll \card{\operatorname{Dom}f}.
\]
Many direct-table inputs then share one output.  The work is to find a small
intermediate value that identifies the right output class.  A small range is the
reason to look for compression; it is not, by itself, a construction or a speed
proof.

The conclusions below are deliberately heuristic and tied to the listed finite
parameters.  This document may be enlarged later when more functions and
ciphertext implementations have been tested.

\section{Notation and what the numbers mean}

We use the notation of the paper.  Thus
\([t]=\{0,1,\ldots,t-1\}\), the usual case is
\[
  f:[t]^\ell\longrightarrow[t],
\]
and the paper searches for unary maps \(p_i:[t]\to\Z\) and an outer map
\(Q\) such that
\[
  b(x_1,\ldots,x_\ell)=\sum_{i=1}^{\ell}p_i(x_i),
  \qquad f(x_1,\ldots,x_\ell)=Q(b(x_1,\ldots,x_\ell)).
\]
Its outer span is
\[
  E=\max_{[t]^\ell}b-\min_{[t]^\ell}b,
\]
so the paper's direct and decomposed record counts are \(t^\ell\) and
\(\ell t+E+1\), respectively.

Some experiments below have unequal input alphabets, for example two binary
strings or \(m\in[0,t-1]\), \(d\in[1,t-1]\).  Only in those cases we use the
more general notation
\(D=D_1\times\cdots\times D_\ell\); the direct count is then \(\card D\).
There is no separate output symbol: ``\#out'' in the tables means the number
of different values of \(f\) on the stated domain.

For a rounded additive construction, the program forms integer unary tables
\(p_i\), enumerates every input, and accepts an outer table only when
\[
 b(x)=b(x')\quad\Longrightarrow\quad f(x)=f(x')
\]
for all tested \(x,x'\).  For comparison networks, histograms, and packed
counts, ``Ours'' is the sum of the consecutive unary-LUT intervals used by the
stated evaluator.  The same LUT may be reused in an implementation, but
repeated evaluations are counted conservatively.  Public additions and
public-constant multiplications are not LUT evaluations.

These are structural record counts, not ciphertext timings.  Parameter growth,
noise, packing, key switching, memory access, and parallelism can change the
practical result.  In both tables, ``Direct'' is the number of records in one
full LUT on the original Cartesian input domain.  ``Gain'' is
\[
  \text{Direct}/\text{Ours}.
\]
``LUT evals.'' is the number of functional-bootstrap LUT evaluations, not
circuit depth: independent evaluations are still counted separately.  In
Table~\ref{tab:elementary}, \(M\) is the integer scale used to round the unary
features.  In Table~\ref{tab:general}, ``Alt.'' is the record count of the
explicit intermediate-value route described below, and ``Alt./ours'' compares
that route with the paper's route.

\section{How the experiments were run}

The deterministic program is
\texttt{expanded\_function\_study/run\_experiments.py}.  It uses only the
Python standard library.  For every row it enumerates the complete finite domain,
evaluates the direct definition and the proposed evaluator, checks every output,
and writes the result to \texttt{results/all\_results.csv} and
\texttt{results/results.json}.

The base run contains 76 parameterized candidates.  Forty-five satisfy its
mechanical rule ``maximum error zero and Direct/Ours at least 2.''  Simple
affine sums followed by clipping or thresholding are still omitted because they
do not need the nonlinear decomposition studied here.

This revision also runs
\texttt{expanded\_function\_study/compare\_intermediate\_routes.py}.  It checks
59 function/parameter rows against an explicit competing route that first forms
an algebraic, sorted, or histogram intermediate and then uses one final LUT.
The competitor is deliberately favoured: its ciphertext multiplications are
reported but assigned zero LUT records.  A row is retained only when every
finite input is correct, Direct/Ours is at least 2, and Ours is strictly smaller
than the competing record count.  The script retains 19 elementary rows and 26
further non-trivial rows; all 14 rejected comparisons remain in the CSV and
JSON output.

For the transcendental ratios, every direct output was computed independently at
80 and 110 decimal digits.  The two integer results had to agree at every input.
Exact integer identities, such as \(\log 9/\log 3=2\), are handled explicitly by
snapping only after agreement to 65 decimal places.  This is a high-precision
finite check, not a symbolic proof about transcendental constants.

\section{Low-depth and deeper elementary compositions}

\subsection{The finite transform pattern}

A useful finite-domain pattern uses logarithms to turn products or ratios into
additive features and then applies an outer reconstruction table.  The same
search pattern is useful for a somewhat deeper expression
\[
  f(x_1,\ldots,x_\ell)
  =F\bigl(\phi_1(x_1)+\cdots+\phi_\ell(x_\ell)\bigr).
\]
For a scale \(M\), the program rounds the unary features to integer tables,
adds them, and checks whether the resulting code separates all output classes.
The accepted outer LUT is defined on the whole interval from the smallest to the
largest code.  Consequently, a function is rejected when the precision needed
for separation makes that interval too large, even if its output range is small.

\subsection{Retained experiments}

Table~\ref{tab:elementary} contains every retained elementary-composition row.
The notation and record count are the same as in Section 2.  In the bivariate
ratio rows \(x,y\in\{1,\ldots,t-1\}\), so the direct count is \((t-1)^2\).

\begin{longtable}{@{}L{0.39\textwidth}rrrrrr@{}}
\caption{Retained elementary compositions; every listed input was checked.}
\label{tab:elementary}\\
\toprule
Function and parameters & Direct & \#out & Ours & Gain & \shortstack{LUT\\evals.} & \(M\)\\
\midrule
\endfirsthead
\toprule
Function and parameters & Direct & \#out & Ours & Gain & \shortstack{LUT\\evals.} & \(M\)\\
\midrule
\endhead
\(\lfloor m/d\rfloor\), \(m\in[0,63]\), \(d\in[1,63]\)
  &4032&64&1192&3.38&3&89\\
\(\lfloor(x_1x_2x_3)^{1/3}\rfloor\), \(x_i\in[1,15]\)
  &3375&15&1252&2.70&4&103\\
\(\lfloor a^{1/b}\rfloor\), \(a,b\in[2,16]\)
  &225&4&106&2.12&3&15\\
\(\lfloor a^{1/b}\rfloor\), \(a,b\in[2,32]\)
  &961&5&234&4.11&3&27\\
\(\lfloor\log_a b\rfloor\), \(a,b\in[2,32]\)
  &961&6&347&2.77&3&61\\
\(\lfloor\log(1+x)/\log(1+y)\rfloor\), \(t=32\)
  &961&6&343&2.80&3&60\\
\(\lfloor\log(1+x^2)/\log(1+y^2)\rfloor\), \(t=32\)
  &961&10&435&2.21&3&56\\
\(\lfloor\sqrt{1+x^2}/\sqrt{1+y^2}\rfloor\), \(t=32\)
  &961&22&437&2.20&3&42\\
\(\lfloor\operatorname{asinh}(x)/\operatorname{asinh}(y)\rfloor\), \(t=32\)
  &961&5&411&2.34&3&78\\
\(\lfloor\sqrt{\log(1+x)/\log(1+y)}\rfloor\), \(t=32\)
  &961&3&339&2.83&3&119\\
\(\lfloor\sqrt[3]{\log(1+x)/\log(1+y)}\rfloor\), \(t=32\)
  &961&2&335&2.87&3&175\\
\(\lfloor\sqrt{\operatorname{asinh}(x)/\operatorname{asinh}(y)}\rfloor\), \(t=32\)
  &961&3&409&2.35&3&155\\
\(\lfloor\log(1+x)/\log(1+y)\rfloor\), \(t=64\)
  &3969&7&913&4.35&3&152\\
\(\lfloor\log(1+x^2)/\log(1+y^2)\rfloor\), \(t=64\)
  &3969&12&1187&3.34&3&148\\
\(\lfloor\sqrt{1+x^2}/\sqrt{1+y^2}\rfloor\), \(t=64\)
  &3969&45&1311&3.03&3&108\\
\(\lfloor\operatorname{asinh}(x)/\operatorname{asinh}(y)\rfloor\), \(t=64\)
  &3969&6&1001&3.97&3&178\\
\(\lfloor\sqrt{\log(1+x)/\log(1+y)}\rfloor\), \(t=64\)
  &3969&3&911&4.36&3&303\\
\(\lfloor\sqrt[3]{\log(1+x)/\log(1+y)}\rfloor\), \(t=64\)
  &3969&2&913&4.35&3&456\\
\(\lfloor\sqrt{\operatorname{asinh}(x)/\operatorname{asinh}(y)}\rfloor\), \(t=64\)
  &3969&3&981&4.05&3&347\\
\bottomrule
\end{longtable}

The intermediate-route screen also keeps all 19 rows.  For example, first
forming \(x_1x_2x_3\) in the geometric-mean row needs two ciphertext
multiplications and leaves a product interval of 3375 values for the final
cube-root/floor LUT.  Even when those two multiplications are treated as free,
the competing LUT has 3375 records, whereas the logarithmic route has 1252.
For the bivariate rows, the conservative competitor is one packed-pair LUT;
the costs of first evaluating the two transforms are omitted.  Thus the
displayed Gain is also a lower bound on the record advantage over that
competitor.

The last fourteen rows use deeper elementary compositions: their unary feature
already contains a logarithm of a logarithm, a square followed by logarithms, or
a hyperbolic inverse followed by a logarithm; three rows then add another square
root or cube root before flooring.  They nevertheless need only two inner unary
tables and one outer unary table on the tested domains.

We also tested the suggested trigonometric ratios rather than assuming that they
would work.  The decompositions were exact, but none reached the 2-fold rule.
For
\[
 \left\lfloor\frac{\sin(\pi x/(2t))}{\sin(\pi y/(2t))}\right\rfloor
\]
the ratios were 0.97, 1.07, and 0.94 at (t=16,32,64).  The corresponding
tangent ratios were 1.13, 1.45, and 0.86, and the cosine ratios were 0.97,
1.07, and 0.94.  Adding an outer square root to the sine ratio still gave only
1.18, 1.07, and 0.94.  These rows remain in the raw result files.  They are a
useful warning that greater composition depth and a small output range do not
guarantee a short consecutive code interval.

\section{Further non-trivial functions}

The complete algorithm and every experiment line are available at
\url{https://github.com/3458463857-tech/A-New-Framework-for-Efficient-Multivariate-Functional-Bootstrapping.git}.
The repository upload itself is left to the authors; the prepared upload folder
contains the exact script, raw outputs, and this English source and PDF.

The comparison here is intentionally harder than Direct versus Ours.  For each
row we wrote down an obvious route that first constructs a useful intermediate
and then applies one final LUT.  Ciphertext multiplications in that route are
listed in the output file but, conservatively, contribute zero LUT records.  We
keep a function only when Ours is still strictly smaller than Alt.  This is a
record-count statement, not a proof of lower latency on every FHE platform.

In the first part of Table~\ref{tab:general},
\(z_{(0)}\le\cdots\le z_{(5)}\) are the sorted inputs for \(t=8,\ell=6\).
The paper's full sorting route uses a verified 12-comparator network, hence
\(12(2t-1)=180\) records.  The competing route rank-packs all
\(\binom{t+\ell-1}{\ell}=1716\) sorted tuples and uses one more LUT, for 1896
records in total.  Range and central-four sum need only a seven-comparison
min/max network; packing the resulting pair gives the smaller competing count
\(105+8^2=169\).  The other Alt. counts are stated exactly in the CSV.

For the second part, \(c_v=\card{\{i:x_i=v\}}\), with \(t=5,\ell=8\).
Building all five counts costs \(8\cdot5^2=200\) records in both routes.  The
competitor then packs the five counts in base 9 and uses one LUT of size
\(9^5=59049\), for 59249 records.  Our route instead applies small local tests
or count comparisons.  Every input tuple was enumerated, and the sorting
networks were also checked on all binary inputs by the zero-one principle.

\begin{longtable}{@{}L{0.36\textwidth}rrrrrr@{}}
\caption{Functions retained after the explicit intermediate-route comparison.}
\label{tab:general}\\
\toprule
Exact function and tested parameters & Direct & \#out & Ours & Alt. & Alt./ours & \shortstack{LUT\\evals.}\\
\midrule
\endfirsthead
\toprule
Exact function and tested parameters & Direct & \#out & Ours & Alt. & Alt./ours & \shortstack{LUT\\evals.}\\
\midrule
\endhead
Maximum \(z_{(5)}\)
 &262144&8&75&262144&3495.25&5\\
Minimum \(z_{(0)}\)
 &262144&8&75&262144&3495.25&5\\
Lower median \(z_{(2)}\)
 &262144&8&180&1896&10.53&12\\
Upper median \(z_{(3)}\)
 &262144&8&180&1896&10.53&12\\
Second smallest \(z_{(1)}\)
 &262144&8&180&1896&10.53&12\\
Second largest \(z_{(4)}\)
 &262144&8&180&1896&10.53&12\\
Range \(z_{(5)}-z_{(0)}\)
 &262144&8&105&169&1.61&7\\
Interquartile range \(z_{(3)}-z_{(1)}\) (fixed finite-sample convention)
 &262144&8&180&1896&10.53&12\\
Top-two sum \(z_{(5)}+z_{(4)}\)
 &262144&15&180&1896&10.53&12\\
Bottom-two sum \(z_{(0)}+z_{(1)}\)
 &262144&15&180&1896&10.53&12\\
Top-three sum \(z_{(3)}+z_{(4)}+z_{(5)}\)
 &262144&22&180&1896&10.53&12\\
Central-four sum \(z_{(1)}+z_{(2)}+z_{(3)}+z_{(4)}\)
 &262144&29&105&169&1.61&7\\
Median-pair gap \(z_{(3)}-z_{(2)}\)
 &262144&8&180&1896&10.53&12\\
Median absolute deviation: lower median of
\(\{|x_i-z_{(2)}|:1\le i\le6\}\)
 &262144&4&450&262414&583.14&30\\
Median-neighborhood count
\(\sum_{i=1}^{6}\mathbf1[|x_i-z_{(2)}|\le1]\)
 &262144&6&270&262414&971.90&18\\
Mode, with the smallest symbol chosen on a tie, \(t=5,\ell=8\)
 &390625&5&268&59249&221.08&44\\
Distinct count \(\sum_{v=0}^{4}\mathbf1[c_v>0]\)
 &390625&5&245&59249&241.83&45\\
Plurality margin: largest \(c_v\) minus second largest \(c_v\)
 &390625&8&353&59249&167.84&49\\
Modal frequency \(\max_v c_v\)
 &390625&7&268&59249&221.08&44\\
Singleton-symbol count \(\sum_v\mathbf1[c_v=1]\)
 &390625&5&245&59249&241.83&45\\
Repeated-symbol count \(\sum_v\mathbf1[c_v\ge2]\)
 &390625&4&245&59249&241.83&45\\
Exact-two-symbol count \(\sum_v\mathbf1[c_v=2]\)
 &390625&5&245&59249&241.83&45\\
Heavy-symbol count \(\sum_v\mathbf1[c_v\ge3]\)
 &390625&3&245&59249&241.83&45\\
Least positive frequency \(\min\{c_v:c_v>0\}\)
 &390625&5&313&59249&189.29&49\\
Occupied-frequency range
\(\max_v c_v-\min\{c_v:c_v>0\}\)
 &390625&7&381&59249&155.51&53\\
Number of modal symbols \(\sum_v\mathbf1[c_v=\max_u c_u]\)
 &390625&4&353&59249&167.84&49\\
\bottomrule
\end{longtable}

Several formerly reported functions fail this stricter test and are therefore
absent from the table.  Integer RMS uses 343 records in our coordinate-square
route, but only 295 in the optimistic ``form the square sum, then bootstrap''
route.  The corresponding variance counts are 12733 and 442.  The Winsorized
mean now includes the previously omitted final division LUT: both descriptions
use 223 records, so there is no strict advantage.  Second moment, collision
count, and Gini bucket use 245 records through count LUTs but only 200 if the
five ciphertext count multiplications are treated as free.  Jaccard, Dice,
cosine, and overlap similarly lose by 113 versus 81, 185 versus 153, and 761
versus 729 for the last two.  The four dynamic-program examples are also removed
because a record-only comparison cannot establish dominance over bit/circuit
implementations.

Every value above is generated by
\path{expanded_function_study/compare_intermediate_routes.py} and appears on
one line of
\path{expanded_function_study/intermediate_routes/intermediate_routes.csv}.
The file also contains
the rejected rows, the exact competing route, its LUT evaluations, and the
number of ciphertext multiplications that the conservative comparison ignored.
This makes the claim deliberately narrow: the retained functions win this
explicit structural record comparison, while ciphertext timings are still
needed before claiming an end-to-end speedup.

\section{What the experiments suggest}

The present data support the following limited conclusions.
\begin{enumerate}
  \item A much smaller range than domain is the common reason compression is
  possible.  A successful evaluator then exposes a short code interval or a
  small sufficient state.
  \item Low composition depth is a good search guide, especially when logarithms
  turn products, ratios, variable roots, or variable logarithms into additive
  features.  It is not a guarantee: the tested sine, cosine, tangent, harmonic
  mean, and two-input geometric-mean cases are concrete counterexamples under
  the stated metric.
  \item The surviving general examples use order statistics or categorical
  histogram tests.  Low-degree moments and normalized set counts are useful
  intermediates, but in this comparison they favour the competing route.
  \item Reachable-state count is not enough.  A sparse encoding can still force a
  large consecutive LUT interval, as the rejected variance experiment shows.
  \item The right conclusion is function-by-function and parameter-by-parameter.
  Structural screening should be followed by a same-platform, correctness-checked
  TFHE implementation before claiming an end-to-end acceleration.
\end{enumerate}

This is an open empirical catalogue.  Later versions may add further elementary
compositions, application-specific statistics, or ciphertext measurements, but
new rows should use the same rule: publish the definition and code, keep negative
results, verify every feasible finite input, and label sampled or model-based
evidence plainly.

\section{Files and reproduction}

From the prepared GitHub folder, run
\begin{center}
\texttt{python expanded\_function\_study/run\_experiments.py}\\
\texttt{python expanded\_function\_study/compare\_intermediate\_routes.py}\\
\texttt{python expanded\_function\_study/verify\_reported\_tables.py}
\end{center}
to regenerate the 76-row base study, the 59 intermediate-route comparisons,
and the source-to-table checks.  The folder also contains the Chapter 4
compression code, its exact representation records and tests, and the TFHE-rs
Table 11 prototype.  The implementation for Section 3 is not included in this
release and will be open-sourced in a later update.

\end{document}
